\chapter{Results and Discussion}

This chapter reports the results gathered so far. The alpha testing results (unit and system testing) and the beta testing results are as of mid-November 2026. The user acceptance testing (UAT) has been carried out, and the results of its preliminary survey and feedback survey are reported below.

\section{AI Mood Performance}
\label{sec:ai-results}
This section reports what was measured for the mood AI, in the order of Table~\ref{tab:ai-tests}. Every figure is reported whether or not it looks good, and none is rounded up. A published figure, where one exists, is a reference point and not a target, because it comes from different data. The methods are in Section~\ref{sec:ai-eval}, and Section~\ref{sec:ai-summary} compares the results with the requirements of Table~\ref{tab:ai-req}.

\subsection{Sound-Type Step}
\label{sec:ai-soundtype}
Table~\ref{tab:ai-sound} gives the average recall of speech, of music, and of neither on the six MUSAN sources that were not used for tuning. The application uses the 5-second window. The published comparison is 93.3\% to 96.6\% on average over four road conditions~\cite{won2017adaptive}, which was measured on other data.

\begin{table}[htbp]
    \centering
    \caption{Sound-type step on sources not used for tuning}\label{tab:ai-sound}
    \small
    \renewcommand{\arraystretch}{1.2}
    \begin{tabular}{|c|c|}
        \hline
        \textbf{Window} & \textbf{Average recall} \\
        \hline
        1 s & 0.959 \\
        \hline
        3 s & 0.977 \\
        \hline
        5 s & 0.972 \\
        \hline
    \end{tabular}
\end{table}

Speech and music were recognized at 0.98 to 0.99. Quiet, the class of neither, was the weakest, at 0.905 at 1 second, 0.95 at 3 seconds, and 0.93 at 5 seconds, and its test is a single source of 87 files. With noise added at 20, 10, 5, and 0~dB, the recall of speech and of music at 5 seconds fell from 0.99 and 0.98 to 0.89 and 0.83 at 0~dB, and the errors mostly became Quiet. MUSAN noise is not the noise of a car, so these figures do not show how the step behaves in a real cabin.

\subsection{Music Head}
\label{sec:ai-music}
Table~\ref{tab:ai-deam} compares the three heads on DEAM at the cutoff of 5.0. The shipped head is the first row.

\begin{table}[htbp]
    \centering
    \caption{Music head on DEAM (10 folds, songs held out)}\label{tab:ai-deam}
    \small
    \renewcommand{\arraystretch}{1.2}
    \begin{tabular}{|>{\raggedright\arraybackslash}p{.42\textwidth}|c|c|c|c|}
        \hline
        \textbf{Head} & \textbf{Spearman} & \textbf{Calm found} & \textbf{Lively found} & \textbf{Balanced} \\
        \hline
        Trained on windows, read on windows (shipped) & 0.790 & 0.832 & 0.804 & 0.818 \\
        \hline
        Trained on whole excerpts, read on the whole excerpt & 0.786 & 0.832 & 0.806 & 0.819 \\
        \hline
        Trained on whole excerpts, read on windows & 0.783 & 0.909 & 0.650 & 0.779 \\
        \hline
    \end{tabular}
\end{table}

The difference in the Spearman correlation at song level (mean and 95\% interval from a bootstrap) is 0.007 [0.000, 0.015] for the shipped head minus the whole-excerpt head read on windows, and 0.004 [$-0.003$, 0.011] for the shipped head minus the whole-excerpt head, so the shipped head is statistically indistinguishable from the whole-excerpt head at song level. At the level of the window, the Spearman correlation is 0.732 for the shipped head and 0.725 for the whole-excerpt head. For the whole-excerpt head, the notebooks also give a CCC of 0.757, an RMSE of 0.81 against 1.29 for always guessing the mean, and an accuracy of 0.820 in separating high from low. CCC and RMSE were not computed for the shipped head. On 58 full-length songs kept out of training, the whole-excerpt head fell to a CCC of 0.613 and an accuracy of 0.67, a small sample of about $\pm$0.12, and the shipped head was not scored on them. Songs by the same artist may sit on both sides of the folds, so these figures may be optimistic. The published references are an accuracy of 81\% to 87\% in separating high from low arousal~\cite{sun2023segmented} and a Spearman correlation of about 0.59 for music to music~\cite{zhang2026generalisation}. The balanced accuracy of 0.818 and the Spearman correlation of 0.790 of the shipped head are not on the same data, so they are reference points and not targets that were met.

On the 767 PMEmo clips, 0.739 of the clips are rated above 0.5. Table~\ref{tab:ai-pmemo} gives the song-level readings.

\begin{table}[htbp]
    \centering
    \caption{Music head on PMEmo (767 chorus clips, never used for training)}\label{tab:ai-pmemo}
    \footnotesize
    \setlength{\tabcolsep}{3pt}
    \renewcommand{\arraystretch}{1.2}
    \begin{tabular}{|>{\raggedright\arraybackslash}p{.26\textwidth}|c|c|c|c|c|c|}
        \hline
        \textbf{Head} & \textbf{Spearman} & \textbf{AUC} & \textbf{Calls Lively} & \textbf{Calm found} & \textbf{Lively found} & \textbf{Balanced} \\
        \hline
        Trained on windows (shipped) & 0.664 & 0.840 & 0.837 & 0.455 & 0.940 & 0.698 \\
        \hline
        Trained on whole excerpts, whole clip & 0.665 & 0.847 & 0.819 & 0.505 & 0.933 & 0.719 \\
        \hline
        Trained on whole excerpts, read on windows & 0.669 & 0.847 & 0.752 & 0.590 & 0.873 & 0.732 \\
        \hline
    \end{tabular}
\end{table}

The Spearman differences all have intervals that include zero: $-0.006$ [$-0.020$, 0.009] for the shipped head minus the whole-excerpt head read on windows, and $-0.002$ [$-0.018$, 0.016] for the shipped head minus the whole-excerpt head. The ranking transfers well, with a Spearman correlation of about 0.66, above the published reference of 0.59 for music to music~\cite{zhang2026generalisation} on different data. The Calm and Lively labels transfer less well. At the cutoff of 5.0 of DEAM, the shipped head calls 83.7\% of the clips Lively, against 73.9\% that were rated above 0.5, and it finds only 45.5\% of the Calm clips.

\textbf{Cutoff of the music head.} The per-second ratings of DEAM put neutral at 3.9 to 4.4 on the scale of the song, against the midpoint of 5.0, and the two labels agree on 69\% of the songs. For the whole-excerpt head, the balanced accuracy stays between 0.77 and 0.82 whatever the cutoff, and only the split between Calm and Lively moves. The value of 5.0 is therefore a working value, and it was checked against car recordings (Section~\ref{sec:ai-rating}).

\subsection{Speech Head}
\label{sec:ai-speech}
Table~\ref{tab:ai-speechres} gives the leave-one-actor-out validation over all 24 RAVDESS actors, with 1,248 clips. The CCC is the mean over the 24 actors, with the standard deviation in parentheses. ``Calm found'' is the share of sad clips read below the cutoff of 0, ``Lively found'' is the share of high-energy clips read above it, and ``Balanced'' is their average. Neutral clips sit exactly on the cutoff, so they are reported separately and are not scored as right or wrong. The x-vector with a linear head scored best on CCC and on balanced accuracy, so it was used. The training run printed a convergence warning for the small networks, because the limit of 500 iterations was reached before they settled, so their rows are a fair but not fully tuned comparison. The published comparison is a CCC of about 0.63 to 0.67 on IEMOCAP~\cite{zhou2023arousal}. The result here is 0.604 on a different data set, which is slightly below that range. The head finds about 43\% of Calm speech, so soft talking can still read Lively.

\begin{table}[htbp]
    \centering
    \caption{Speech head, three embeddings with a linear head and with a small network (leave one actor out)}\label{tab:ai-speechres}
    \footnotesize
    \setlength{\tabcolsep}{3pt}
    \renewcommand{\arraystretch}{1.2}
    \begin{tabular}{|>{\raggedright\arraybackslash}p{.26\textwidth}|c|c|c|c|c|}
        \hline
        \textbf{Embedding and head} & \textbf{CCC (SD)} & \textbf{Calm found} & \textbf{Lively found} & \textbf{Balanced} & \textbf{Neutral read below 0} \\
        \hline
        x-vector, linear head (used) & 0.604 (0.139) & 0.432 & 0.970 & 0.701 & 0.365 \\
        \hline
        x-vector, small network & 0.293 (0.193) & 0.484 & 0.802 & 0.643 & 0.438 \\
        \hline
        ECAPA, linear head & 0.545 (0.136) & 0.365 & 0.972 & 0.668 & 0.260 \\
        \hline
        ECAPA, small network & 0.241 (0.176) & 0.458 & 0.759 & 0.609 & 0.385 \\
        \hline
        YAMNet mean and maximum, linear head & 0.454 (0.137) & 0.109 & 0.983 & 0.546 & 0.146 \\
        \hline
        YAMNet mean and maximum, small network & 0.356 (0.102) & 0.411 & 0.878 & 0.645 & 0.375 \\
        \hline
    \end{tabular}
\end{table}

\textbf{Stability on windows.} The speech head was trained on whole clips, so it was checked on 150 MUSAN speech clips of 10 seconds by comparing the reading of a window of 5.295 seconds with the reading of the whole clip. The spread (standard deviation) of the readings of the whole clips is 0.265. For the first window, the correlation is 0.924, the mean absolute difference is 0.095, and the label of Calm or Lively agrees in 92.7\% of the clips. For the last window, the correlation is 0.883, the difference is 0.100, and the label agrees in 93.3\%. For the mean of the two windows, the correlation is 0.988, the difference is 0.032, and the label agrees in 96.7\%. The first and the last window agree with each other on the label in 87.3\% of the clips. A single window is therefore a little less stable than the whole clip, which is one reason for the confirm rule. This measures stability only and not accuracy.

\subsection{Mixed Audio}
\label{sec:ai-mixed}
The 150 pairs of speech and music were mixed with the music from 12~LU above the speech to 36.4~dB below it.
\begin{itemize}
    \item \textbf{What the sound-type step hears.} It hears no music at 36.4~dB below the speech, hears it in 24\% of the windows at 22.7~dB below, and reports both sources in 50\% to 77\% of the windows between 18 and 7~LU below.
    \item \textbf{How one source changes the label of the other.} Adding music changes the label of the speech head in 9\% to 25\% of the cases between 22.7~LU below and an equal level. Adding speech changes the label of the music head in 28\% to 41\% of the cases between 22.7 and 7~LU below, and it lowers the reading by about half a rating point, so the speech reading is the steadier one.
    \item \textbf{The two orders.} The speech-first order and the music-first order give different labels in 71\% to 75\% of the mixtures in that band. Speech first keeps the clean label of the speech in 81\% to 91\% of them, and music first keeps the clean label of the music in 59\% to 72\%.
\end{itemize}
The music-head figures were measured with the shipped head, whose reading is taken on the first and last window of 5.295 seconds of each clip of 10 seconds and averaged. With RMS in place of LU, the label changes are 30\% to 43\% and the Lively share is 48\%, so the finding does not depend on the measure of loudness. The speech head calls 85\% of the MUSAN speech Lively, and the music head calls 48\% of the music Lively, so one cutoff on two data sets may not mean the same arousal. MUSAN speech is read speech with no pauses.

\subsection{Confirm Rule and Update Interval}
\label{sec:ai-confirm}
The simulation gave $N=2$ as the smallest number of identical labels in a row that meets the target of 0.10 false flips for each minute. It draws from a saved pool of the pipeline's own labels, built from 200 speech clips and 250 music clips of 10 seconds (160,000 samples at 16~kHz) from MUSAN. From each clip, two windows of 5.295 seconds were taken, the first and the last, which overlap by about 0.6 seconds and so are nearly independent but not fully. Each window ran through the full pipeline with the music cutoff of 5.0, and the category of a clip came from its clean reading, the average of the readings of its two windows. The Quiet pool is 200 windows of synthetic pink noise at about $-50$~dBFS, a level the proponents chose. The speech comes from librivox (100) and us-gov (100), and the music from fma, fma-western-art, hd-classical, jamendo, and rfm (50 each), each a random segment with a fixed random seed and an RMS above 0.01. These are the same MUSAN sources that were used to tune and test the sound-type cutoffs, so the clip sources overlap with that work. Table~\ref{tab:ai-pool} gives the share of windows that received each label.

\begin{table}[htbp]
    \centering
    \caption{Share of windows given each label (the input of the simulation)}\label{tab:ai-pool}
    \small
    \renewcommand{\arraystretch}{1.2}
    \begin{tabular}{|l|c|c|c|}
        \hline
        \textbf{Category of the clip (windows)} & \textbf{Lively} & \textbf{Calm} & \textbf{Quiet} \\
        \hline
        Speech Lively (352) & 0.95 & 0.04 & 0.01 \\
        \hline
        Speech Calm (48) & 0.21 & 0.79 & 0.00 \\
        \hline
        Music Lively (232) & 0.95 & 0.05 & 0.00 \\
        \hline
        Music Calm (268) & 0.05 & 0.94 & 0.01 \\
        \hline
        Quiet (200) & 0.00 & 0.00 & 1.00 \\
        \hline
    \end{tabular}
\end{table}

The category comes from the reading of the same pipeline, so the table shows how consistent single windows are with their own clip and not how well they agree with human raters. The Quiet result is a test on synthetic pink noise and not on car noise, and the group of Calm speech has only 48 windows.

Table~\ref{tab:ai-sim} gives the result for updates every 5.295 seconds. At $N=1$ the false flips for each minute were above the target, and $N=2$ was the first value to meet it with independent errors and 60-second scenes.

\begin{table}[htbp]
    \centering
    \caption{Simulated stored label, updates every 5.295 seconds (300 timelines for each row)}\label{tab:ai-sim}
    \footnotesize
    \setlength{\tabcolsep}{3pt}
    \renewcommand{\arraystretch}{1.2}
    \begin{tabular}{|c|l|c|c|c|c|c|}
        \hline
        \textbf{Scene length} & \textbf{Errors} & \textbf{$N$} & \textbf{Accuracy} & \textbf{Lag (s)} & \textbf{False flips per minute} & \textbf{Reversals per hour} \\
        \hline
        60 s & independent & 1 & 0.897 & 3.396 & 0.671 & 82.09 \\
        \hline
        60 s & independent & 2 & 0.869 & 9.185 & 0.071 & 18.00 \\
        \hline
        60 s & by clip & 1 & 0.896 & 3.576 & 0.422 & 55.93 \\
        \hline
        60 s & by clip & 2 & 0.834 & 9.784 & 0.130 & 21.03 \\
        \hline
        25 s & independent & 2 & 0.730 & 8.704 & 0.067 & 50.28 \\
        \hline
        25 s & by clip & 2 & 0.709 & 8.950 & 0.127 & 50.82 \\
        \hline
        120 s & independent & 2 & 0.923 & 9.370 & 0.087 & 9.88 \\
        \hline
        120 s & by clip & 2 & 0.877 & 10.356 & 0.142 & 14.75 \\
        \hline
    \end{tabular}
\end{table}

\textbf{What $N=2$ costs and what it does not fix.} The rule cuts the false flips about tenfold for independent errors (from 0.671 to 0.071 for each minute at 60 seconds), but it adds about 5.8 seconds of lag (from 3.396 to 9.185 seconds) and lowers the accuracy by about three points (from 0.897 to 0.869). In the worst case (by clip), the false flips stay at 0.127 to 0.142 for each minute, above the target of 0.10, because a persistent error is not a flicker and no value of $N$ fixes it. A larger $N$ lowers the false flips further but raises the lag: at $N=6$ with 60-second scenes the lag is 26.7 seconds by clip and 27.7 seconds independent.

\textbf{Update interval.} An update every 10 seconds is worse at the same $N$: at $N=2$ with independent 60-second scenes, the accuracy is 0.785 and the lag 15.593 seconds, against 0.869 and 9.185 seconds at 5.295 seconds. In the separate grid, the accuracy was highest at 5.295 seconds for every scene length and fell with each longer interval, as Table~\ref{tab:ai-grid} shows. The bootstrap test found every longer interval clearly worse than the best, so the rule of the largest interval that is not clearly worse gave 5.295 seconds in all three cases. The tolerances of 0.02 and 0.05 in accuracy that were examined are the proponents' own. The energy column of the grid was left empty, because no measurement of energy on a phone was made. This result supports the window of 5.295 seconds.

\begin{table}[htbp]
    \centering
    \caption{Accuracy of the stored label by update interval (grid of the simulation)}\label{tab:ai-grid}
    \small
    \renewcommand{\arraystretch}{1.2}
    \begin{tabular}{|l|c|c|c|c|}
        \hline
        \textbf{Scene length} & \textbf{5.295 s} & \textbf{10 s} & \textbf{20 s} & \textbf{60 s} \\
        \hline
        25 s & 0.856 & 0.792 & 0.673 & 0.496 \\
        \hline
        60 s & 0.898 & 0.871 & 0.814 & 0.640 \\
        \hline
        120 s (smoothing span of 15 s at 5.295 s) & 0.923 & 0.899 & 0.870 & 0.770 \\
        \hline
    \end{tabular}
\end{table}

\subsection{Models on the Phones}
\label{sec:ai-phones}
Table~\ref{tab:ai-phoneres} gives the checks of the models on the phones. These are timings of debug builds on one Android phone and one iPhone and not benchmarks.

\begin{table}[htbp]
    \centering
    \caption{Checks of the models on the phones}\label{tab:ai-phoneres}
    \small
    \renewcommand{\arraystretch}{1.2}
    \begin{tabular}{|>{\raggedright\arraybackslash}p{.34\textwidth}|>{\raggedright\arraybackslash}p{.56\textwidth}|}
        \hline
        \textbf{Check} & \textbf{Result} \\
        \hline
        Converted models match the originals (final package, 10 golden clips) & YAMNet: largest difference 0.000006 in any score and 0.000040 in any embedding. X-vector: largest embedding difference 0.004689, on values of a typical size of about 8 to 12. Largest change in the reading of the speech head: 0.00003. \\
        \hline
        Golden clips reproduce the notebook on both phones & 10 of 10 clips, with a tolerance of 0.001 \\
        \hline
        Android, golden test, time for each window & median 289 ms \\
        \hline
        iPhone, golden test, time for each window & median 52 ms, worst 285 ms (the first, warm-up window) \\
        \hline
        Android, live use (debug build) & The sound-type step takes about 120 to 140 ms for each window, a speech window adds about 160 to 190 ms, and there was no visible stutter. \\
        \hline
        Live sanity check on a phone & In a quiet room, music and speech each gave the expected label. \\
        \hline
    \end{tabular}
\end{table}

In a live run on an iPhone on the driver side with music playing near the phone, every window was scored in 56 to 70 ms in total. The quiet windows had speech scores of 0.006 to 0.012 and music scores of 0.002, below both cutoffs, so the raw label was Quiet. The music windows had music scores of 0.74 to 0.92 and speech scores near 0, so the music head ran. Its readings of 4.148 and 4.665 gave Calm, and its readings of 5.176 and 5.352 gave Lively, on either side of the cutoff of 5.0. The stored label changed only after two identical raw labels in a row, and single raw labels left it unchanged. This is one short run, shown as evidence that the chain works on the phone and not as a result of accuracy.

The mood lock was checked on one iPhone. The lock held after the mood drifted to Quiet before the pickup, and cancelling every accepted request restarted the microphone. It was not checked on Android. On the iPhone, the music that is playing pauses for less than a second when the microphone starts and when it stops, and then resumes by itself. Five off and on switches were tested, and each gave a short pause followed by an automatic resume.

\subsection{Researcher Rating Check on Car Recordings}
\label{sec:ai-rating}
The two blind raters agreed on 25 of the 30 recordings (Cohen's kappa 0.713), and the label that both gave was taken as the agreed label. The sound-type step told the kind of sound correctly in 29 of the 30 clips. The labels of all 30 clips are listed in Appendix~\ref{app:clips}. Table~\ref{tab:ai-raters} compares the label of the AI with the agreed label.

\begin{table}[htbp]
    \centering
    \caption{Label of the AI against the label that both blind raters gave}\label{tab:ai-raters}
    \small
    \renewcommand{\arraystretch}{1.2}
    \begin{tabular}{|l|c|c|}
        \hline
        \textbf{Kind of sound} & \textbf{Clips with an agreed label} & \textbf{AI matched the agreed label} \\
        \hline
        Quiet & 3 & 3 \\
        \hline
        Music & 12 & 5 \\
        \hline
        Speech & 10 & 1 \\
        \hline
        All & 25 & 9 (36\%) \\
        \hline
    \end{tabular}
\end{table}

The AI matched the agreed label on 9 of 25 clips (36\%), with a kappa of $-0.07$. The AI called all 11 speech clips Lively, and the raters called 9 of the 10 speech clips on which they agreed Calm. The kappa between the researcher who recorded the clips, who was not blind, and each blind rater was 0.714 and 0.767, and the kappa of the three raters together (Fleiss) was 0.731. These figures are supplementary.

The result does not support the Lively and Calm labels of the AI on these recordings. The cutoffs (speech 0.15, music 0.20, music arousal 5.0, and speech 0) were not changed after it, and it is reported as it is. It agrees with two limits known before the check. First, the speech head was trained on acted speech and finds only about 43\% of Calm speech (Section~\ref{sec:ai-speech}). Second, the cutoff of 0 for speech is the neutral point, so soft talking that is not sad can read Lively. As an exploratory look only, the readings of the music head ranked the 12 agreed music clips in the order of the raters (AUC 0.88), but the cutoff of 5.0 sat above the readings of 7 of the 8 clips that the raters called Lively. For speech, only one agreed clip was Lively, so the cutoff for Lively speech could not be tested. With 25 agreed clips and two blind raters from the research team, this is a small check. It does not show that the labels match the way students rate sound in a car, and the research does not claim agreement with human ratings.

\subsection{Matching and Privacy Tests}
\label{sec:ai-matchtests}
The conditions of Table~\ref{tab:ai-accept} run in the matchRide Cloud Function, which is covered by the Cloud Functions test suite and nine verification scripts (Section~\ref{sec:phases}). The matchRide script runs 19 checks and all passed: 5 for the mood filter, 3 for the vehicle filter, and 11 for the other conditions. The mood checks cover a passenger who accepts only Calm and is refused a Lively driver, a passenger who accepts Calm and Quiet and is matched with a Quiet driver and refused a Lively driver, a driver whose label is unknown who is accepted (fail open), and a passenger with no preference who accepts any driver. The Flutter test suite has 76 tests of the mood code (Table~\ref{tab:unit}). In the system test cases (Section~\ref{sec:system-results} and Appendix~\ref{app:systemcases}), S-C5 covers the filter in the whole application: a passenger who accepts only Quiet is tested with a Lively driver, a Quiet driver, and a driver with no mood label. S-F3 checks that audio is sampled only while the driver is available and that only the mood label is stored.

\subsection{Summary of Results Against the Requirements}
\label{sec:ai-summary}
Table~\ref{tab:ai-summary} compares the results with the requirements of Table~\ref{tab:ai-req}. A requirement is marked ``Not supported'' when the evidence does not show it, and ``Partly met'' when part of it is shown and part is not.

\begin{table}[htbp]
    \centering
    \caption{Results of the mood AI against its requirements}\label{tab:ai-summary}
    \small
    \renewcommand{\arraystretch}{1.2}
    \begin{tabular}{|c|>{\raggedright\arraybackslash}p{.60\textwidth}|>{\raggedright\arraybackslash}p{.18\textwidth}|}
        \hline
        \textbf{No.} & \textbf{Result} & \textbf{Status} \\
        \hline
        R1 & Sound-type recall 0.972 at 5 s. Music head Spearman 0.790 on DEAM and 0.664 on PMEmo. Speech head CCC 0.604. On 30 car recordings, the AI matched the agreed label on 9 of 25 clips (kappa $-0.07$). & Not supported on car recordings \\
        \hline
        R2 & Speech and music are rated on the same axis. The two heads sit on different baselines (85\% of speech and 48\% of music read Lively on MUSAN). & Partly met \\
        \hline
        R3 & Converted models match the originals, and 10 of 10 golden clips are reproduced within 0.001. Median time for a window is 289 ms on Android and 52 ms on iPhone. & Met \\
        \hline
        R4 & Audio is analyzed on the phone and only the label is stored. 76 Flutter tests of the mood code passed. Motorcycles never run the microphone. & Met \\
        \hline
        R5 & $N=2$ gives 0.071 false flips for each minute at a lag of 9.185 s for independent errors. For persistent errors the false flips stay at 0.127 to 0.142 for each minute. & Partly met \\
        \hline
        R6 & The matchRide script passes all 19 checks, including the fail-open cases. & Met \\
        \hline
    \end{tabular}
\end{table}

\subsection{Limitations of the Mood AI Results}
\label{sec:ai-limits}
Each limitation below is a fact about what was and was not tested.

\textbf{Validity of the labels}
\begin{itemize}
    \item \textbf{Not validated against ratings in a car.} The tests do not show that the labels match how students rate the sound in a car, or that the cutoffs suit podcasts in particular. The rating check did not support the Lively and Calm labels (Section~\ref{sec:ai-rating}). It used two blind research-team members and the recorder, who was not blind, which is few. All clips were played through one car stereo and recorded in one car, with one iPhone 16 in one position, so the results say nothing about other cars, phones, or positions. The cutoffs are working values.
    \item \textbf{Disagreement among raters caps the accuracy.} Human raters agree only moderately on arousal (alpha of 0.28 to 0.66 in DEAM~\cite{aljanaki2017deam}), so no accuracy above that ceiling is claimed.
    \item \textbf{The two heads sit on different baselines.} The speech head calls 85\% of the MUSAN speech Lively and the music head calls 48\% of the music Lively, so one cutoff on two data sets may not mean the same arousal. The Calm and Lively labels of the music head transfer less well than its ranking (Section~\ref{sec:ai-music}).
    \item \textbf{The passenger needs behind the three moods are assumptions.} No survey of music preferences was conducted, and whether mood matching improves rides was not measured.
\end{itemize}

\textbf{Data and training}
\begin{itemize}
    \item \textbf{The speech head was trained on acted speech.} RAVDESS has two fixed sentences read by actors, not podcasts. The head finds about 43\% of Calm speech, and soft talking can still read Lively. The calm clips of RAVDESS were left out because no published anchor exists for them.
    \item \textbf{Quiet was tested on synthetic noise.} MUSAN noise and added noise are not the noise of a car engine, and a very noisy car can read as Quiet, because the errors under heavy noise fall to Quiet. The test of Quiet is one source of 87 files.
    \item \textbf{The cutoffs of the sound-type step were tuned on MUSAN and not on car audio.} The audibility thresholds for background music were measured on television and not in a car.
    \item \textbf{The rule for mixed audio is the proponents' own.} Speech first has indirect support only. Between 22.7 and 7~LU below the speech, 71\% to 75\% of the mixtures receive a different label under music first.
\end{itemize}

\textbf{Behavior of the label}
\begin{itemize}
    \item \textbf{The first label can be wrong.} It is taken from one window, so one odd window at the start can set it. It corrects after two agreeing windows, but in the worst case of the simulation a persistent error is not corrected by the confirm rule.
    \item \textbf{The microphone hears conversation in the car.} Only the mood word leaves the phone, but the talk of the passengers and the driver is analyzed as speech. No rule separates media from conversation, and both are rated as speech. Separating them is left for future work, and no method has been named.
    \item \textbf{There is no logic for the boundaries of songs.} A song changes section about every 25 seconds on average~\cite{lyu2026sections}, but the application does not look for boundaries.
    \item \textbf{Mood is a fail-open filter.} A wrong but known label can keep a passenger who chose a specific mood from seeing that driver.
    \item \textbf{The matcher does not check how old the label is,} although the label has an update time.
    \item \textbf{The label belongs to the account and not to the phone.} Availability belongs to the account of the driver, so two phones that are signed in to one account both start the microphone.
\end{itemize}

\textbf{Platform and measurement}
\begin{itemize}
    \item \textbf{Battery and energy were not measured.} The team makes no energy claim. The update interval was chosen from accuracy and lag only.
    \item \textbf{The timings are from debug builds} on one Android phone and one iPhone.
    \item \textbf{The mood lock was checked on one iPhone only.} It was not checked on Android.
    \item \textbf{On iPhone, the music pauses briefly when the microphone starts and stops.} Starting the microphone switches the phone to recording mode, which iOS does for any application that records. The pause lasts less than a second and the music then resumes by itself. Releasing the audio session when the microphone stops was added and tested, but the pause is reduced and not removed.
    \item \textbf{The application needs iOS 16 or later.} iPhone testers install through free Apple IDs, so push notifications could not be delivered on iPhone in testing (Section~\ref{sec:scope}).
\end{itemize}

\FloatBarrier
\section{Reading the Test Results}
\label{sec:readresults}
The tests reported in the rest of this chapter are the unit tests, the system test cases, and the beta tester tasks. Each result is recorded under one fixed label, and a label means the same thing wherever it appears. Table~\ref{tab:resultlabels} states the meanings once, so that the tables that follow can be read without further explanation.

\begin{table}[H]
    \centering
    \caption{Meaning of the Result Labels}
    \label{tab:resultlabels}
    \small
    \renewcommand{\arraystretch}{1.25}
    \begin{tabular}{|>{\raggedright\arraybackslash}p{0.15\textwidth}|>{\raggedright\arraybackslash}p{0.19\textwidth}|>{\raggedright\arraybackslash}p{0.50\textwidth}|}
        \hline
        \textbf{Label} & \textbf{Used for} & \textbf{Meaning} \\
        \hline
        Pass & Unit tests, system cases, beta tasks & The actual result matched the expected result written before the test. A system case on a physical device is a Pass only when all five rounds succeeded. \\
        \hline
        Fail & Unit tests, system cases, beta tasks & The actual result differed from the expected result. Every Fail is entered in the defect log with a severity (Section~\ref{sec:phases}); a beta tester also writes what was seen instead. \\
        \hline
        Skipped & Unit tests & The test did not run because an input it needs was not supplied. It is neither a Pass nor a Fail. \\
        \hline
        Could not try & Beta tasks (optional) & The tester did not do one of the optional tasks (O1 to O3). It is neither a Pass nor a Fail. \\
        \hline
        N/A & System cases, per device & The case does not apply to that device, or cannot be tested on it because of a limit of the test setup and not because of how the application behaved. It is not a Fail. \\
        \hline
    \end{tabular}
\end{table}

Pass rates count only what was actually run. A Pass counts in favor, a Fail counts against, and an N/A or Could not try result is left out of the count, which is why the iPhone columns of Table~\ref{tab:system} show 28 cases run instead of 32, and why the beta pass rate ignores Could not try. These results are still reported separately, so that a reader can see how many cases were not run and why, and none of them is hidden inside a pass rate. In this study, N/A applies to the four notification cases (S-E1 to S-E4) on the two iPhones, because the free Apple IDs used could not receive push notifications (Section~\ref{sec:scope}). No unit test was skipped in this study.

A Pass means that the application behaved as the expected result said, in the setting where the case was run, which is an emulator or a physical device. It does not prove that no defect exists, and a result on an emulator is not a result on a physical device. For this reason, the emulator and device results are reported separately.

\section{Alpha Testing: Unit Testing}
\label{sec:unit-results}
Unit testing was run with \texttt{flutter test} and the Cloud Functions test suite. All 390 Flutter tests passed on the final build, with no failures and none skipped, which is a 100\% pass rate. All nine Cloud Functions verification scripts were run on the final build and all passed (matchRide, geofence, Places, Ride Tokens, Routes API, alternative drop-off ranking, the Firestore and Storage security rules, and the fuel price parser). Table~\ref{tab:unit} shows the Flutter tests by area.

\begin{table}[htbp]
    \centering
    \caption{Flutter Unit Tests by Area}
    \label{tab:unit}
    \small
    \renewcommand{\arraystretch}{1.2}
    \begin{tabular}{|l|c|c|}
        \hline
        \textbf{Area} & \textbf{Total cases} & \textbf{Result} \\
        \hline
        Accounts and registration & 46 & Pass \\
        Ride flow and matching & 93 & Pass \\
        Tokens, fuel, and CES period & 54 & Pass \\
        Admin, reports, and ratings & 70 & Pass \\
        Mood detection (mood code only, without ride matching) & 76 & Pass \\
        Shared UI and safety & 51 & Pass \\
        \hline
    \end{tabular}
\end{table}

\section{Alpha Testing: System Testing}
\label{sec:system-results}
System testing ran 32 cases. On the emulator, all 32 passed on the first run, with none failed (a 100\% first-run pass rate). On the three physical devices (a Samsung Galaxy A06 5G, an iPhone 17, and an iPhone 15 Pro Max), each case was run in five rounds and passed only when all five rounds succeeded. All 32 cases passed on the Android phone. The iPhones could not receive push notifications, because the proponents used free Apple IDs (Section~\ref{sec:scope}), so the four notification cases (S-E1 to S-E4) could not be tested on them and are marked N/A. The other 28 cases passed on both iPhones, with none failed. Three cases (S-B5, S-B7, and S-D2) were run for one round only. Table~\ref{tab:system} shows the results by area.

\begin{table}[htbp]
    \centering
    \caption{System Testing Results by Area}
    \label{tab:system}
    \small
    \renewcommand{\arraystretch}{1.2}
    \begin{tabular}{|l|c|>{\centering\arraybackslash}p{1.6cm}|>{\centering\arraybackslash}p{1.6cm}|>{\centering\arraybackslash}p{1.6cm}|>{\centering\arraybackslash}p{1.6cm}|}
        \hline
        \textbf{Area} & \textbf{Cases} & \textbf{Emulator passed} & \textbf{Android passed} & \textbf{iOS 1 passed} & \textbf{iOS 2 passed} \\
        \hline
        B1 Accounts and admin & 6 & 6 & 6 & 6 & 6 \\
        B2 Ride flow & 8 & 8 & 8 & 8 & 8 \\
        B3 Matching filters & 6 & 6 & 6 & 6 & 6 \\
        B4 Tokens and fuel & 4 & 4 & 4 & 4 & 4 \\
        B5 Notifications and chat & 4 & 4 & 4 & N/A & N/A \\
        B6 Safety and privacy & 4 & 4 & 4 & 4 & 4 \\
        \hline
        \textbf{Total} & \textbf{32} & \textbf{32} & \textbf{32} & \textbf{28 run} & \textbf{28 run} \\
        \hline
    \end{tabular}
    \par\smallskip{\footnotesize N/A: not tested, because iOS devices could not receive push notifications (free Apple IDs). Chat itself works on iOS inside the application, but a message is not announced when the application is closed or in the background. iOS 1: iPhone 17. iOS 2: iPhone 15 Pro Max.}
\end{table}

Table~\ref{tab:measure} gives the measurements recorded during system testing. Requests reached the driver in about 1.6 to 1.8 seconds and the passenger's screen updated about 1.2 to 1.6 seconds after the driver accepted. The physical devices were faster than the emulator in every timing measurement. No crashes or freezes were recorded on any device.

\begin{table}[htbp]
    \centering
    \caption{Measurements Recorded During System Testing}
    \label{tab:measure}
    \small
    \renewcommand{\arraystretch}{1.2}
    \begin{tabular}{|>{\raggedright\arraybackslash}p{0.34\textwidth}|c|c|c|c|}
        \hline
        \textbf{Measurement} & \textbf{Emulator} & \textbf{Android} & \textbf{iOS 1} & \textbf{iOS 2} \\
        \hline
        Request to driver prompt (median of 5 runs) & 1.8 s & 1.7 s & 1.7 s & 1.6 s \\
        Accept to passenger update (median of 5 runs) & 1.6 s & 1.4 s & 1.2 s & 1.3 s \\
        Cold start to home screen (median of 3 runs) & 2.7 s & 0.9 s & 1.0 s & 0.9 s \\
        Crashes and freezes & 0 & 0 & 0 & 0 \\
        App size (phone storage) & -- & 199 MB & 197 MB & 197 MB \\
        \hline
    \end{tabular}
    \par\smallskip{\footnotesize Emulator: Google Pixel 8 Pro. Android: Samsung Galaxy A06 5G. iOS 1: iPhone 17. iOS 2: iPhone 15 Pro Max. App size was read from each phone's storage settings.}
\end{table}

The Android app was installed from Android Studio as a debug build. A debug build is not size-optimized, and its package includes native libraries for three CPU architectures (arm64-v8a, armeabi-v7a, and x86\_64) so that one build runs on different Android phones. The debug APK was 109 MB, and the 199 MB shown in the phone's settings also includes the unpacked libraries, compiled code, and app data. The iOS build contains only the architecture of the device it is installed on. For comparison, a release build with one APK per CPU architecture is 32.0 MB (armeabi-v7a), 36.6 MB (arm64-v8a), and 39.9 MB (x86\_64), so the Android figure reported in testing is larger than what a user would install from a release. \TODO{Update the debug APK size and the release APK sizes for the final build, which includes the larger AI models.}

\subsection{Limits Found}
The main limit found in alpha testing was the difficulty of performing matched rides with physical devices. Matching needs a driver and a passenger device in the right places at the same time. Three cases that need such rides (S-B5, S-B7, and S-D2) were run for one round only, against five rounds for the other cases. Each round of these cases needed a driver to drive a vehicle with a passenger device in the matching area, and the proponents limited it to one round to save fuel and time.

\section{Beta Testing}
\label{sec:beta-results}
Twelve volunteers were recruited and all twelve finished the feedback survey. Each tester used one physical device with SanCarGo installed, taken from the ten devices in Table~\ref{tab:betadevices}, and two Android emulators, one playing the passenger and one the driver. The physical device was used for the tasks that did not involve inbound or outbound trips, and the emulators were used for the trips. The iPhones could not show push notifications (Section~\ref{sec:scope}), so the testers tried the notification and chat tasks on the emulators, which can receive them. Only four iPhones were available, so the iOS testers shared devices. The testers signed up themselves, using dummy accounts in the USC email domain (usc.edu.ph) whose credentials the proponents gave them. The average session lasted 40 minutes, against the 45 minutes planned. Eleven of the twelve testers (91.7\%) never use ridesharing applications, and one (8.3\%) sometimes does, so most ratings come from people new to this kind of application.

\subsection{Task Results}
Overall, all of the task results recorded as Pass or Fail were Pass (201 of 201, or 100\%), above the 80\% target of Table~\ref{tab:betapass}. All of the required tasks (T1 to T3 and W1 to W11) were reported as Pass by every tester. Of the optional tasks (O1 to O3), 91.7\% were reported as Pass and 8.3\% as Could not try (three task results in total), which are left out of the pass rate as Table~\ref{tab:resultlabels} states. No crashes were reported, so the condition of no crash on the core tasks was met.

\subsection{Experience Ratings}
Testers rated seven statements about their experience from 1 (strongly disagree) to 5 (strongly agree). Table~\ref{tab:beta-ratings} gives the averages. All were above 4. The highest was the understanding of why location and microphone access are requested (4.6), and the lowest was the readability of the screens and the map (4.1).

\begin{table}[htbp]
    \centering
    \caption{Beta Testing Experience Ratings (Average of 12 Testers)}
    \label{tab:beta-ratings}
    \small
    \renewcommand{\arraystretch}{1.2}
    \begin{tabular}{|l|c|}
        \hline
        \textbf{Statement} & \textbf{Average (1 to 5)} \\
        \hline
        The app responded quickly & 4.3 \\
        Notifications arrived when I expected them & 4.5 \\
        The error messages told me what to do next & 4.4 \\
        The screens and the map were easy to read & 4.1 \\
        The emergency help button was easy to find & 4.5 \\
        I understood why the app asked for location and microphone access & 4.6 \\
        I felt my personal information was handled safely & 4.5 \\
        \hline
    \end{tabular}
\end{table}

\subsection{Open Responses}
All testers answered that the application did not crash, freeze, or lose their place, and all answered that they would use SanCarGo to get to or from campus, giving reasons such as saving money, smart transportation for students, free rides, less traffic, possible environmental benefits, and a new way of meeting people. The parts testers liked most were the interface, the rating feature, accurate location tracking, the role tutorials, the dynamic mood, the geofence, the responsive notifications, and the chat. The most confusing parts were what to expect from the mood feature, some interface text, delayed notifications, the lack of route options on outbound trips, mood buttons that glitched, and delayed screens after completing a ride. Testers suggested more pick-up and drop-off points, better mood classification, a dark mode, live sharing, a calling feature, and different route options for outbound trips.

\subsection{Pre-UAT Dry Run}
Towards the end of beta testing, three people whom the proponents reached out to personally, such as family members, tried the application as a dry run ahead of the UAT (Section~\ref{sec:phases}). It was not a usability evaluation and it was not technical.

\TODO{Add what the three people did in the dry run and what the proponents saw: the part of the flow they tried, any problem found, and any change made.}

\subsection{Limits Found}
Two limits were found in beta testing: visual interface issues on Android devices, and differing opinions on mood classification, which appeared both among the most confusing parts and among the suggested improvements. In addition, the rides were played on emulators with simulated GPS, so the results say nothing about how real GPS drifts on a moving vehicle, battery use and heat, or whether live tracking survives a phone in a pocket. Twelve testers is a small group, and most had not used a ridesharing application before.

\TODO{Add the result of the team checks and the severity of the problems reported by testers, which are not in the results received so far.}


\section{User Acceptance Testing}
\label{sec:uat-results}
\subsection{Preliminary Survey}
\label{sec:prelim-results}
The preliminary survey was answered by 98 of the 378 enrolled Computer Engineering students (25.9\%), which is 51.3\% of the target of 191 respondents (Section~\ref{sec:prelim}). With 98 respondents, the margin of error at a 95\% confidence level is about 8.5\%, assuming $p=0.5$ and the finite population correction, and not the planned 5\%. The results describe the respondents and not all the students of the department. Table~\ref{tab:prelim-results} shows who answered and who chose to be a tester.

\begin{table}[htbp]
    \centering
    \caption{Respondents of the preliminary survey and the respondents who chose to be testers}\label{tab:prelim-results}
    \small
    \renewcommand{\arraystretch}{1.15}
    \begin{tabular}{|>{\raggedright\arraybackslash}p{0.34\textwidth}|>{\centering\arraybackslash}p{0.26\textwidth}|>{\centering\arraybackslash}p{0.26\textwidth}|}
        \hline
        \textbf{Characteristic} & \textbf{All respondents ($n=98$)} & \textbf{Chose to be testers ($n=46$)} \\
        \hline
        \multicolumn{3}{|l|}{\textit{Year level}} \\
        \hline
        Year 1 & 27 (27.6\%) & 16 (34.8\%) \\
        \hline
        Year 2 & 31 (31.6\%) & 20 (43.5\%) \\
        \hline
        Year 3 & 22 (22.4\%) & 10 (21.7\%) \\
        \hline
        Year 4 & 18 (18.4\%) & 0 (0.0\%) \\
        \hline
        \multicolumn{3}{|l|}{\textit{Place of residence}} \\
        \hline
        Cebu City & 30 (30.6\%) & 15 (32.6\%) \\
        \hline
        Lapu-Lapu City & 19 (19.4\%) & 10 (21.7\%) \\
        \hline
        Mandaue City & 19 (19.4\%) & 10 (21.7\%) \\
        \hline
        Liloan & 12 (12.2\%) & 4 (8.7\%) \\
        \hline
        Talisay City & 9 (9.2\%) & 4 (8.7\%) \\
        \hline
        Consolacion & 9 (9.2\%) & 3 (6.5\%) \\
        \hline
        \multicolumn{3}{|l|}{\textit{Operating system of the phone}} \\
        \hline
        Android & 47 (48.0\%) & 20 (43.5\%) \\
        \hline
        iOS (iPhone) & 37 (37.8\%) & 20 (43.5\%) \\
        \hline
        Both & 14 (14.3\%) & 6 (13.0\%) \\
        \hline
        \multicolumn{3}{|l|}{\textit{Mobile network carrier}} \\
        \hline
        Globe or GOMO & 56 (57.1\%) & 25 (54.3\%) \\
        \hline
        Smart, TNT, or Sun & 34 (34.7\%) & 17 (37.0\%) \\
        \hline
        DITO & 8 (8.2\%) & 4 (8.7\%) \\
        \hline
        \multicolumn{3}{|l|}{\textit{Role chosen (testers only)}} \\
        \hline
        Passenger only & -- & 35 (76.1\%) \\
        \hline
        Driver and passenger & -- & 11 (23.9\%) \\
        \hline
        \multicolumn{3}{|l|}{\textit{Vehicle of the drivers (drivers only, $n=11$)}} \\
        \hline
        Car, diesel & -- & 4 \\
        \hline
        Car, gasoline & -- & 2 \\
        \hline
        Motorcycle, gasoline & -- & 5 \\
        \hline
    \end{tabular}
\end{table}

Forty-six of the 98 respondents (46.9\%) chose to be testers, and 52 did not. The testers came from Years 1 to 3: none of the 18 respondents of Year 4 chose to be a tester. Of the testers, 35 chose to be passengers only and 11 chose to be both drivers and passengers, which gives about one driver for every three passengers (11 to 35). Of the 11 drivers, 6 drove cars (4 diesel and 2 gasoline) and 5 rode motorcycles (all gasoline), and 4 of the 5 motorcycle drivers had a spare helmet to lend to a passenger. The phones of the testers were evenly divided between the two platforms: 20 Android phones, 20 iPhones, and 6 respondents who had both, so the application had to work on both.

The respondents lived mostly in Cebu City (30), Lapu-Lapu City (19), and Mandaue City (19), and most used Globe or GOMO (56) or Smart, TNT, or Sun (34) for mobile data. The landmarks that they named near the start of their daily commute were mostly malls and public markets, and the most frequent were SM Consolacion (5 respondents), E-Plaza in Liloan (4), and Pacific Mall (3). One respondent left the landmark blank.

\subsection{Participation and Conduct}
\label{sec:uat-conduct}
Forty-six students took part in the UAT, against the target of about 48 (25\% of the 191 target respondents, Section~\ref{sec:phases}). The target was already a realistic expectation and not a statistical sample, and the number reached was close to it.

The UAT departed from the plan in two ways. First, it was planned for two weeks but was held for one week. The review of the REC took longer than the proponents expected, and the test had to be compressed into the time that remained. Second, only a few participants were willing to take part in actual rides. The controlled setting confined the rides to the campus routes of Table~\ref{tab:rules-by-setting}, and the participants did not find it practical to drive around the campus: the routes were out of the way for them and would take some time, since, as Computer Engineering students, they stayed mostly at the Bunzel building where their classes were held. The participants were willing to use the application and to test it, in particular its user interface, and most of this use was done while stationary, but they were not willing to make actual rides around the campus.

These departures limit what the UAT can show. The results describe the usability and the perception of the application by the participants, and they say little about how the application behaves on actual rides, which the system tests on physical devices and the beta tests on emulators covered only in part. \TODO{Add the number of rides that were matched and completed during the UAT, and the completion and failure rates.}

\subsection{Feedback Survey}
\label{sec:uat-feedback-results}
All 46 participants answered the feedback survey. Thirty-nine (84.8\%) never use ridesharing applications and 7 (15.2\%) sometimes do, so most ratings come from people who were new to this kind of application, as in the beta testing.

\textbf{System Usability Scale.} A SUS score was computed for all 46 participants. The mean SUS score was 80.4 (standard deviation 4.8, median 80.0), against the target of 68 or higher (Section~\ref{sec:survey}). The lowest score was 70.0 and the highest was 90.0, so all 46 scores were above the usual benchmark of about 68 for average usability. Table~\ref{tab:uat-sus} gives the mean answer to each statement. The statements about difficulty were answered low and those about ease were answered high: the participants disagreed that the application was complex (1.67), that they needed technical support (1.61), and that it was inconsistent (1.41), and they agreed that it was easy to use (4.26) and that they felt confident (4.43). The lowest agreement was with statement 7, that most people would learn to use it very quickly (3.91).

\begin{table}[htbp]
    \centering
    \caption{Mean answer to each System Usability Scale statement (1 = strongly disagree, 5 = strongly agree)}\label{tab:uat-sus}
    \small
    \renewcommand{\arraystretch}{1.15}
    \begin{tabular}{|>{\raggedright\arraybackslash}p{0.62\textwidth}|c|c|}
        \hline
        \textbf{Statement} & \textbf{Answers} & \textbf{Mean} \\
        \hline
        1. I would like to use SanCarGo frequently. & 46 & 4.28 \\
        \hline
        2. I found SanCarGo unnecessarily complex. & 46 & 1.67 \\
        \hline
        3. I thought SanCarGo was easy to use. & 46 & 4.26 \\
        \hline
        4. I would need the support of a technical person to use SanCarGo. & 46 & 1.61 \\
        \hline
        5. The various functions in SanCarGo were well integrated. & 46 & 4.39 \\
        \hline
        6. There was too much inconsistency in SanCarGo. & 46 & 1.41 \\
        \hline
        7. Most people would learn to use SanCarGo very quickly. & 46 & 3.91 \\
        \hline
        8. I found SanCarGo very cumbersome to use. & 46 & 2.02 \\
        \hline
        9. I felt very confident using SanCarGo. & 46 & 4.43 \\
        \hline
        10. I needed to learn a lot of things before I could get going with SanCarGo. & 46 & 2.41 \\
        \hline
    \end{tabular}
\end{table}

\textbf{Experience ratings.} Table~\ref{tab:uat-exp} gives the averages of the experience statements, rated from 1 to 5. All were above 3.8. The highest was that the error messages told the participant what to do next (4.57), followed by the understanding of why location and microphone access were requested (4.54) and the feeling that personal information was handled safely (4.52). The lowest was that the emergency help button was easy to find (3.85), the only average below 4.

\begin{table}[htbp]
    \centering
    \caption{UAT experience ratings (average of 46 participants)}\label{tab:uat-exp}
    \small
    \renewcommand{\arraystretch}{1.15}
    \begin{tabular}{|>{\raggedright\arraybackslash}p{0.72\textwidth}|c|}
        \hline
        \textbf{Statement} & \textbf{Average (1 to 5)} \\
        \hline
        The app responded quickly & 4.30 \\
        \hline
        Notifications arrived when I expected them & 4.48 \\
        \hline
        The error messages told me what to do next & 4.57 \\
        \hline
        The screens and the map were easy to read & 4.39 \\
        \hline
        The emergency help button was easy to find & 3.85 \\
        \hline
        I understood why the app asked for location access (and the microphone, for the driver) & 4.54 \\
        \hline
        I felt my personal information was handled safely & 4.52 \\
        \hline
    \end{tabular}
\end{table}

\textbf{Open responses.} All 46 participants answered that the application did not crash, freeze, or lose their place. Thirty-five participants (76.1\%) answered that they would use SanCarGo to get to or from campus and 11 (23.9\%) answered maybe; none answered no. The reasons for maybe were that it depends on how many drivers are available or using the application, on the class schedule, on whether it covers the area where the participant lives, on more pickup points, and on the application staying reliable.

Thirty-five participants named nothing confusing. The others named what Ride Tokens are, finding the token history, how the mood chip works, what the status messages meant, where to start a new request, knowing when the driver accepted, waiting after sending a request, typing a destination against dropping a pin, unclear error messages, and, in two answers, the emulator (setting its location, and which emulator was the driver).

What the participants liked most was the ease and clarity of the flow (13 answers, including the welcome tutorial, clear buttons, and clear trip steps), the look and readability of the screens (6), speed, notifications, and real-time updates (7), the map and the driver information such as name, plate, rating, and chat (5), the safety features such as the emergency button (3), the Ride Token and CES rewards (3), that no payment is involved (2), that the application is made for USC students (2), its reliability (3), the mood detection (1), and the destination search (1).

Twenty-four participants suggested no change. The suggestions of the others were more pickup and drop-off areas or a search for places (5), an arrival time or a countdown for the driver (3), visual changes such as brighter colors, a dark mode, and larger text (4), a tutorial or walkthrough for new users (2), more mood options (2), clearer error messages, an easier way to find the chat, a sound for new requests, a display of the number of passengers waiting, staying logged in longer, and a quick help button (one answer each).

\section{Discussion}
The alpha testing results show that the core logic of the application works as designed. The automated checks of ride matching and the Ride Token system (Objectives 2 and 3) all passed, and the whole-system cases (Objective 1) all passed on the emulator and on the Android phone, and on the two iPhones all the cases that could be run passed (the four notification cases could not be tested on iOS), with the matching, token, notification, and safety flows behaving as specified on the devices that could receive notifications. The response times of about one to two seconds for requests and acceptances are short enough for a ride-matching interaction.

Read against the test levels of ISO/IEC/IEEE 29119 used as the reference framework in Section~\ref{sec:phases}, the evidence is uneven. The unit level is complete: all automated checks passed. The system level is complete on the emulator and on the Android phone, and complete on the two iPhones except for the four notification cases, with each case recorded as a step, an expected result, and an actual result per device, as the standard's documentation practice suggests. Beta testing has produced results from 12 outside volunteers, but the acceptance level, which is the one that tests the application with its intended users, has no results yet. The mood AI was evaluated on held-out data, on mixed audio, on the phones, and in a rating check on car recordings, which is the kind of evidence the standard's guidance on AI-based systems (Part 11) points to for a component that cannot be verified by fixed expected outputs alone~\cite{iso29119site}. The testing was not assessed for conformance with the standard, so these statements describe how the evidence is organized and do not claim certification.

The beta results are positive. All required tasks passed for every tester, every task result that was recorded as Pass or Fail was a Pass, no crashes were reported, and the experience ratings were all above 4 out of 5, which suggests that testers who had mostly never used a ridesharing application could complete the main ride flow. The main weakness they pointed to was the mood feature: its classification, what to expect from it, and its buttons were mentioned among both the most confusing parts and the improvements wanted, which agrees with the rating check of Section~\ref{sec:ai-rating}, where the Lively and Calm labels of the AI did not match the labels of the raters. The models themselves ran as intended on the phones: the golden clips were reproduced on one Android phone and one iPhone (Section~\ref{sec:ai-phones}).

The results also show what is not yet proven. Three system cases were run for one round only. The beta rides were played on emulators with simulated GPS, and the group was small, so the results do not show how the application behaves with real GPS on a moving vehicle. Finally, Objective 4, the evaluation with Computer Engineering students, is met in part by the UAT. The mean SUS score of 80.4 was above the target of 68, no participant reported a crash, and every experience average was above 3.8, with the emergency button the least easy to find (3.85). The UAT, however, reached 46 students, was held for one week and not the planned two, and included only a few actual rides (Section~\ref{sec:uat-conduct}), and most participants had never used a ridesharing application. The results therefore say more about the usability of the application than about its behavior on rides, and the narrow spread of the SUS scores (standard deviation 4.8) comes from a single group of students of one department. \TODO{Rewrite this discussion after the UAT and survey results are in, and relate the findings to the related literature in Chapter 2.}
